% Part VI -- Comparative synthesis, limitations and future directions (target: about 20 pages)
\clearpage
\part{Comparative Synthesis, Limitations and Future Directions}
\label{part:synthesis}
\partbudget{20}{comparison framework and future-directions outline drafted; conclusions await Parts IV--V}

\section{Monolithic, Chiplet-Based and Vertically Stacked Designs Compared}
\label{sec:comparison}

\Cref{tab:three-way} compares the three organisations against the bottleneck--benefit--trade-off test of \cref{sec:rq}. Entries are qualitative and will be revised against the evidence gathered in Part~\ref{part:case-study}.

\begin{table}[htbp]
  \centering
  \caption{Qualitative comparison of the three design approaches (provisional)}
  \label{tab:three-way}
  \small
  \begin{tabularx}{\textwidth}{@{}>{\raggedright\arraybackslash}p{2.6cm}>{\raggedright\arraybackslash}X>{\raggedright\arraybackslash}X>{\raggedright\arraybackslash}X@{}}
    \toprule
    & \textbf{Monolithic die} & \textbf{Chiplets (2D package)} & \textbf{Chiplets with stacked cache} \\
    \midrule
    Main limit addressed & --- & Die size, yield, uneven node scaling & Planar cost and wire length of a large LLC \\
    Communication & All on-die & Die-to-die links through the package & Adds dense vertical links between cache and compute dies \\
    Memory latency & Shortest on-die paths & Extra die-to-die hop to memory & Fewer trips to memory for cache-sensitive work \\
    Cost and yield & Falls steeply with die area & Smaller dies, mixed nodes, tested before assembly & Extra die and bonding step; capacity only where sold \\
    Thermal & One heat source per die & Heat spread over several dies & Stack alters heat path from cores \\
    New trade-off & Scaling ceiling & Latency, link power, packaging & Heat, test complexity, workload dependence \\
    \bottomrule
  \end{tabularx}
\end{table}

\section{Interactions Between Chiplet Organisation and Cache Design}
\label{sec:interactions}
\plannedcontent{Develop the argument that chiplets and stacked cache are complementary: the CCD's small, self-contained die with its own L3 makes it a natural base for a stacked cache, and stacking lets a product line offer large-cache variants without changing the base CCD or the package layout. Discuss the per-CCD L3 boundary: stacked capacity helps threads on that CCD only, which interacts with scheduling and with cross-CCD communication (\cref{sec:d2d}).}

\section{Workload-Dependent Trade-Offs and Remaining Bottlenecks}
\label{sec:remaining-bottlenecks}
\plannedcontent{Synthesise Parts~III--V: which workloads favour more cores, which favour more cache, which favour frequency; the remaining bottlenecks (memory bandwidth, die-to-die latency, heat removal from stacked dies).}

\section{Future Directions in Heterogeneous Integration and Advanced Packaging}
\label{sec:future}

AMD's Hot Chips~33 presentation sketches integration schemes that go beyond stacking cache on cores, including memory stacked on compute, cores stacked on cores or on uncore logic, and splitting or ``folding'' circuits across dies; AMD presents these as schemes that advanced packaging could enable but monolithic design could not \cite{swaminathanAdvancedPackagingEnabling2021}. These are vendor projections, not products.

\plannedcontent{Discuss the architectural questions these raise (heat density of logic-on-logic stacks, testability, design tools, standard die-to-die interfaces). Candidate sources to be identified and read; no claims about specific future products without primary sources.}

\section{Limitations of This Study}
\label{sec:limitations}
\begin{itemize}
  \item The quantitative analysis relies on published measurements and is limited by their experimental conditions and by the availability of comparable data.
  \item Much of the detailed technical information about AMD designs comes from AMD itself; such claims are attributed to AMD and treated as engineering claims rather than independent results.
  \item The study works at the architectural level and does not model physical design, packaging processes or thermal behaviour in detail.
  \item \todo{Add limitations arising from the final evidence set.}
\end{itemize}

\section{Conclusions}
\label{sec:conclusions}
\plannedcontent{Answer RQ1--RQ4 in turn. Write only after Parts IV and V are complete.}

\section{Open Research Questions}
\label{sec:open-questions}
\begin{itemize}
  \item When is the extra complexity of 3D integration worth its reduction in footprint and its denser interconnect, compared with a side-by-side cache chiplet?
  \item How should operating systems and runtimes place threads and data when L3 capacity differs between compute dies?
  \item What limits the stacking of logic on logic, as opposed to SRAM on logic, from a thermal and power-delivery point of view?
  \item \todo{Refine after the evidence review.}
\end{itemize}
